% =====================================================================
% Exemple minimal — TD (classe article)
% Étape 1 : squelette qui compile. Le contenu viendra à l'étape 3.
% Compilation : latexmk -pdf main.tex
% =====================================================================
\documentclass[11pt]{article}

\newcommand{\relativePath}{../..}
\usepackage[nocorrection, relativePath=\relativePath]{\relativePath/td}

\title{Exemple de TD}
\shorttitle{Exemple de TD}
\numero{TD n\textsuperscript{o}~1}
\date{2025--2026}
\discipline{Nom de la matière}
\promotion{Établissement --- Promotion}

\begin{document}

\maketitle

\begin{myexercise}

    \begin{myquestion}
        Énoncé de la question.
    \end{myquestion}

    \begin{mycorrection}
        Corrigé, masqué par l'option \texttt{nocorrection}.
    \end{mycorrection}

\end{myexercise}

% <<< étape 3 : sous-questions, problèmes, rappels, suppléments, options >>>

\end{document}
